\documentclass[12pt]{article}
\usepackage[utf8]{inputenc}
\usepackage[T1]{fontenc}
\usepackage{hyperref}
\title{Soli Language Resources for K-12 Curriculum Development}
\author{Compiled by Claude Sonnet 5.5 (claude-sonnet-5-5) for Andr\'as Kornai}
\date{}
\begin{document}
\maketitle

\section{Introduction}
Soli (ciSoli, ISO 639-3 sby) is a Bantu language of Zambia. It is spoken around Lusaka and in Central Province, near the Tonga and Lenje, and is described as closely related to Tonga. Speaker counts differ by source: 34,000 (Wikipedia, 2010) and 77,000 (Joshua Project, citing Ethnologue 2016). The searches found a New Testament but no school materials or dictionary. This dossier is short.

\section{Primary Educational Resources}
No primers, textbooks or school programmes for Soli were found.

\section{Linguistic and Reference Materials}
\begin{itemize}
\item Wikipedia, Soli language: \url{https://en.wikipedia.org/wiki/Soli_language}
\item Wikipedia, Soli people: \url{https://en.wikipedia.org/wiki/Soli_people}
\end{itemize}
No dictionary or grammar turned up.

\section{Religious and Community Resources}
\begin{itemize}
\item Soli New Testament on eBible.org (CC BY 4.0, copyright 2018 The Word for the World International; formats include PDF, ePub and Sword): \url{https://ebible.org/bible/details.php?id=sby}
\item About page for the Soli New Testament: \url{https://ebible.org/study/content/texts/sby/about.html}
\item Full-text PDF of the New Testament: \url{https://eBible.org/pdf/sby/sby_nt.pdf}
\item Seed Company story on the Soli translation (New Testament finished, Old Testament begun): \url{https://seedcompany.com/stories/a-growing-appetite-for-gods-word}
\item Joshua Project profile: \url{https://joshuaproject.net/people_groups/14977}
\end{itemize}

\section{Digital and Audio Resources}
Wikitongues has a video of a speaker, ``Mary speaking Soli'': \url{https://wikitongues.org/videos/mary_20191012_sby/}. Wikitongues also says some literary standardization exists but is not widespread.

\section{Curriculum Development Context}
The New Testament is the only substantial text in Soli found, and it can supply a corpus and a provisional orthography. Everything else would be new work: the orthography should be checked with the community, and the relationship to Tonga (and so the choice of school language) settled. Zambian language-in-education policy was not researched here.
\end{document}